\section{Sensitivity, strengths and weaknesses, conclusions}\label{sec:eval}

\subsection{Sensitivity analysis}

A result that flips under a plausible change of a nuisance parameter is not a result. Table~\ref{tab:params} in Appendix~\ref{app:params} records, for every parameter, whether it is assumed, estimated or calibrated, where it comes from, and whether it was swept; we report here what the sweeps changed.

\textbf{Curve number.} The capacity ratio of 5.6 is the observed peak hour divided by the design intensity and is therefore independent of CN by construction; it cannot be used as evidence of robustness. What does depend on CN is the volume. Runoff from the three-day event ranges from 2.72 to $3.52\times10^{8}$\,m$^3$ across CN\,75--95, and the time needed to drain that runoff at 36\,mm\,h$^{-1}$ from 6.4 to 8.3\,h, so the statement that the network was overloaded for most of a working day survives the whole bracket. In Task~5 the same bracket moves Haikou's best uniform ratio from 0.155 (CN\,75) to 0.274 (CN\,95): the level moves, the verdict does not.

\textbf{Shape parameter and estimation method.} Zhengzhou's tail is not identified, and the windows we use do not even agree on its sign, so every window is named. On the 1961--2020 daily maxima the MLE gives $\xi=-0.026$ and the L-moment fit $-0.001$; on 1961--2024, $-0.007$ and $+0.030$; on the 1961--2010 training window of Task~3, $-0.042$ and $-0.022$. The spread between estimators inside one window is as large as the spread between windows, which is the honest summary of a 50--60 year record. The 100-year level is therefore quoted as [139.2, 242.5]\,mm, and the 449\,mm total only as an order of magnitude ($1.6\times10^{8}$\,yr, interval 678\,yr to infinity), never as a number a planner should use. The estimator moves the Task~5 design depth by 0.5\% (102.9\,mm by L-moments against 102.4\,mm by MLE on 1961--2020) and Haikou's ratio from 0.181 to 0.179, so no conclusion rests on it.

\textbf{Prediction target.} The most influential choice in Task~3 is the predictand. Quantile versus expected maximum shifts RMSE by 9.1\,mm with a paired bootstrap interval of [3.3, 13.6]\,mm and bias by 21\,mm, more than any distributional assumption; the gain of the trend correction over the stationary expected maximum, 1.7\,mm with interval $[-1.9, 7.0]$, is inside the noise. Both targets are reported and the expected maximum is recommended for planning.

\textbf{Storm-type thresholds.} The type test uses a rarity ratio of 3 and a concentration gap of 0.10, while the observed values are 14.1 and 0.134. Relaxing the thresholds to 10 and 0.05 leaves both criteria violated, so the ``different type'' verdict is robust to the threshold choice.

\textbf{Investment parameters.} The damage exponent and the discount rate drive the Task~5 levels. Varying $\gamma$ over 1.2--1.8 and the discount rate over 2--6\% moves the best uniform ratio into 0.06--0.44 for all three cities (Guangzhou highest, Zhengzhou lowest), so no parameter combination on that range brings uniform capacity to break-even. Using MLE rather than L-moments for the two-year design depth moves the ratio by less than 0.003, so the tail estimator is not what drives the verdict. Anchoring the damage function on the station total (313.4\,mm) instead of the official areal 449\,mm rescales every ratio by 2.465, which lifts Guangzhou's hotspot ratio to 10.0 while leaving the uniform menu at or below 0.69 --- the hotspot conclusion is the robust one.

\subsection{Strengths}

{\small
\begin{itemize}[leftmargin=1.4em,itemsep=1pt,topsep=2pt]
\item \textbf{Reproducibility.} Every number is generated by the scripts in Appendix~\ref{app:code} from archived raw files with fixed seeds; the pipeline decodes both provider conventions and validates them against each other.
\item \textbf{Honest inference.} False-discovery control over one explicitly declared family of 146 tests, bootstrap intervals on the two comparisons that decide the paper, a heterogeneity test before any regional pooling, and an explicit refusal to interpret return periods beyond what the record supports.
\item \textbf{A real forecast experiment.} Task~3 is evaluated on held-out years with identical folds for ten methods, and the comparison identifies \emph{why} methods differ (target definition, bias, ranking) rather than only which wins.
\item \textbf{A decision-relevant result.} Hazard, exposure, vulnerability and cost share units, so the Task~5 policy follows from the Task~1--4 measurements; protecting 2\% of the area reaches 4.05 at a fraction of the cost of a uniform upgrade that reaches 0.28.
\end{itemize}}

\subsection{Weaknesses and how they would be fixed}

{\small
\begin{itemize}[leftmargin=1.4em,itemsep=1pt,topsep=2pt]
\item \textbf{Daily resolution.} The archives are daily, so the hourly peak enters only through the official record; this is why the daily screen ranks Zhengzhou low rather than first. Sub-hourly telemetry would fix it.
\item \textbf{Decomposition is an identity.} The three factors in Task~4 exhaust the loss gap by construction and do not establish causation; only the exposure mechanism is corroborated externally.
\item \textbf{Short held-out window.} The 300-valid-day rule leaves 5--7 test years per city, so RMSE carries a large sampling error and close calls (GEV versus Gumbel) should not be over-read.
\item \textbf{Loss granularity and hydraulics.} Losses are provincial and municipal rather than district-level, and the damage function is calibrated on one event; we compute volumes and capacity ratios, not street depths. A district inventory over several events plus a coupled 1D--2D model would let the exponent be estimated rather than assumed.
\item \textbf{Stationarity.} The hazard model is stationary with an optional linear trend and uses no climate projections; defensible for the baseline given the weak national trends, but a 30-year horizon should carry a scenario range, since observed intensification of sub-daily extremes \cite{ipcc2021} points the same way as our Zhengzhou result.
\end{itemize}}

\subsection{Conclusions}

\textbf{Task 1.} Zhengzhou's rainfall is not increasing in total; it is becoming more concentrated. Wet days fall by 2.0 per decade ($p=0.021$) while daily intensity rises, and no change point is detectable. The July 2021 daily total (188.7\,mm) was a 54-year empirical event and did not break the 1978 record, but the three-day total (313.4\,mm) was the largest in 71 years and the official 449\,mm lies beyond reliable extrapolation. The waterlogging was a capacity failure in time: the peak hour exceeded the one-year design intensity 5.6-fold, leaving $1.4\times10^{8}$\,m$^3$ above capacity, and the three-day runoff needed 7.4\,h to drain.

\textbf{Task 2.} Across 25 cities only 2 annual-total and 1 one-day-extreme trend survive false-discovery control, against 7 intensity and 9 wet-day trends: redistribution, not more rain. Six belts differ in the sign of the total trend (South $+21.5$, Southwest $-27.1$\,mm per decade) while every belt loses wet days, and two clusters capture the split (silhouette 0.336).

\textbf{Task 3.} The predictand matters more than the distribution: forecasting the expected maximum of $n$ years rather than the $T$-year quantile cuts RMSE from 64.8 to 55.7\,mm and bias from $-41.5$ to $-20.5$\,mm. Extreme-value models dominate machine learning on ranking (top-five overlap 4/5 versus 3/5 at best; Spearman 0.75 versus 0.64), while regression removes the bias. Five cities broke their record in the held-out window, none in 2021, which exposes the blind spot of daily indices for sub-daily disasters. Fuzhou, Haikou, Guangzhou and Hohhot carry the highest ten-year exceedance probabilities; Zhengzhou ranks sixteenth on probability and tenth on the composite.

\textbf{Task 4.} The two storms are different types. Zhengzhou is an intensity-type event (1-day return period 134.5 years, 52\% of the event total in one day, 201.9\,mm in one hour); Taiyuan is a persistence-type event (9.5-year return period, concentration at the 6.5th percentile of its own climatology). Losses differ by 23.9, of which 1.07 is the area ratio; the remaining 22.4 splits into population density (2.68), affected share (2.95) and loss per affected person (2.84), and a Shapley decomposition assigns $31.3$, $33.6$ and $32.7\%$ to the three, so removing any one still leaves 8.1--8.9. Rainfall depth differs by only 2.2, so the gap is exposure and vulnerability, not meteorology. Zhengzhou should be screened on hourly intensity, Taiyuan on multi-day accumulation.

\textbf{Task 5.} Uniform capacity is never the efficient first purchase once damage is priced per unit of overflow volume: across the 378-point menu the best ratio is 0.28 (Guangzhou), 0.18 (Haikou) and 0.10 (Zhengzhou), and no social plan exceeds 0.26. Protecting the 2\% most exposed area, where damage density is twenty times the average, yields 4.05, 2.62 and 1.50 for the same money --- the reverse of the hazard ordering, because absolute damage density decides whether a hotspot investment pays. Under a CVaR$_{95}$ cap no city invests without one and every cap is met below 0.28, so a cap is a constraint to satisfy at least cost, and the least-cost instrument is the hotspot. The plan is therefore three-speed: hotspots and closure protocols within five years, source control where redevelopment makes it cheap within fifteen, and large conveyance only where exposure is extreme. Because a daily screen cannot see a Zhengzhou-type event, every city should archive sub-hourly rainfall.